% PHYS 4604 Research Proposal Template
% Modeled on the NSF GRFP Graduate Research Plan Statement.
%
% NSF formatting rules (checked October 2, 2026; see
% https://www.nsf.gov/funding/initiatives/grfp/how-to-apply):
%   - 2 pages maximum, INCLUDING references and figures
%   - 8.5 x 11 in. paper, PDF only
%   - 1 in. margins with nothing in them (no header, footer, name, or page number)
%   - Computer Modern or Times New Roman, 11 pt or larger
%   - At least single spacing
%   - Separate "Intellectual Merit" and "Broader Impacts" headings
%
% Do not shrink \linespread, use negative \vspace, or drop below 11 pt to fit.
% Cut words instead.

\documentclass[11pt,letterpaper]{article}
\usepackage[margin=1in]{geometry}
\usepackage{amsmath}
\usepackage{graphicx}
\usepackage[compact]{titlesec}
\usepackage[hidelinks]{hyperref}

\titleformat*{\section}{\normalsize\bfseries}
\setlength{\parindent}{0pt}
\setlength{\parskip}{0.5em}

\pagestyle{empty} % NSF: no page numbers or other text in the margins

\begin{document}

\begin{center}
{\large\bfseries Your Title Names the Question, Not the Subfield}
\end{center}

\section*{Motivation and Background}
Why this problem, and why now? Start with why it matters, in words a physicist
outside your subfield would follow. Then name the specific gap in what is known,
and cite the work that shows it is a gap~\cite{example2024}. Keep this to about
half a page; it is the first thing to cut if you run long.

\section*{Research Question}
\textbf{State your research question in one or two sentences, in bold, so a
reviewer can find it in seconds.} Say what you would measure, calculate, or
build, and what each possible answer would mean.

\section*{Approach and Methods}
How would you actually do it? Name the data, instrument, model, or calculation.
Explain why this method can answer the question before explaining how it works.
Show the scope is feasible for an early graduate student with a mentor. This
section gets the most space, about three quarters of a page.

% Optional: one figure, only if it replaces a paragraph.
% \begin{figure}[h]
%   \centering
%   \includegraphics[width=0.6\textwidth]{figure.pdf}
%   \caption{One sentence on what the reader should see.}
% \end{figure}

\section*{Intellectual Merit}
What will the field learn if this works, and if it does not? Connect the result
back to the gap you named above. Then show you can do it: the research
experience, skills, and methods you already have that prepare you for this
project. NSF rates the applicant's potential, not only the project's.

\section*{Broader Impacts}
Who benefits beyond your subfield, how many, and how would you know it worked?
Build on outreach, teaching, mentoring, or open tools you have already done,
rather than promising something new.

\begin{thebibliography}{9}
\setlength{\itemsep}{0pt}
\bibitem{example2024}
A.~Author and B.~Author, ``Title of the paper,''
\textit{Phys.\ Rev.\ Lett.} \textbf{132}, 000000 (2024).
\end{thebibliography}

\end{document}
